% ============================================================
%  Background section: PyTorch 2.0 Compiler Stack + Triton CPU
%  Two-column IEEE conference format
% ============================================================
\documentclass[conference]{IEEEtran}

% ---- packages -----------------------------------------------
\usepackage[T1]{fontenc}
\usepackage[utf8]{inputenc}
\usepackage{microtype}
\usepackage{cite}
\usepackage{amsmath}
\usepackage{graphicx}
\usepackage{booktabs}
\usepackage{url}
\usepackage{xcolor}
\usepackage{enumitem}
\usepackage{balance}

\begin{document}

% ============================================================
\section{Background}
\label{sec:background}
% ============================================================

% ------------------------------------------------------------
\subsection{The PyTorch~2.0 Compiler Stack with Triton CPU}
\label{sec:bg_pytorch}
% ------------------------------------------------------------

PyTorch 2.0 introduced a production-grade compilation pathway that
transforms eager-mode model definitions into high-performance native
code without requiring any changes to model source~\cite{ansel2024pytorch2}.
The stack pairs \emph{TorchDynamo}—a Python bytecode–level graph
capture frontend—with \emph{TorchInductor}, a backend that schedules
and generates kernels for the captured computation graph.
By default, Inductor targets CPU workloads by emitting C++ loop nests
that are compiled and executed directly.
In this work we replace the C++ CPU codegen backend with the
\emph{Triton CPU} backend~\cite{tillet2019triton, triton_cpu_repo},
motivated by two properties that are central to our approach:

\begin{enumerate}[leftmargin=*, label=\arabic*., itemsep=2pt]

    \item \textbf{Inductor-integrated extensible codegen.}
          Triton CPU is an experimental backend that plugs into
          Inductor's codegen abstraction layer.  Unlike the C++ backend,
          it exposes a well-defined, stage-by-stage lowering pipeline
          into which new target-specific transformations can be inserted
          without modifying core Inductor code.  This makes it possible
          to augment the generated kernel call-sites with the per-socket
          memory placement and metadata annotations required for
          tensor-parallel execution.

    \item \textbf{MLIR-based compiler and runtime extensibility.}
          Triton CPU is built on the MLIR compiler
          infrastructure~\cite{lattner2021mlir}, which organises
          compilation as a cascade of modular dialect lowerings.
          Because both the compiler passes and the runtime launcher
          are written against well-specified MLIR and C++ interfaces,
          new execution models—such as the NUMA-aware thread pool
          introduced in this work—can be grafted onto the runtime
          without invasive changes to the Triton compiler itself.

\end{enumerate}

The following subsections describe the end-to-end compilation pipeline,
from model definition to kernel execution on CPU, focusing on the
unmodified upstream behaviour.  Points where this work extends the
pipeline are noted inline and elaborated in Section~\ref{sec:approach}.

% ............................................................
\subsubsection{TorchDynamo: Graph Capture}
\label{sec:bg_dynamo}
% ............................................................

TorchDynamo operates by intercepting Python bytecode at the interpreter
level before it is executed~\cite{ansel2024pytorch2}.
When a model is decorated with \texttt{torch.compile}, Dynamo installs
a frame evaluation hook that observes each Python instruction.
As the model's forward pass runs, Dynamo records the sequence of ATen
tensor operations into an acyclic dataflow graph—the \emph{FX graph}—in
which each node represents a primitive operator and edges carry tensor
values between producers and consumers.
Control flow that depends on runtime tensor values, such as
data-dependent branches, is handled by splitting the graph at the
branch point and re-tracing the relevant sub-graph when the condition
changes.  Shapes and data types are propagated through the graph by
executing each node symbolically against abstract tensor
(\emph{FakeTensor}) proxies, so that downstream stages can reason about
memory layouts and sizes without running real computation.

% ............................................................
\subsubsection{TorchInductor: Lowering and Optimization}
\label{sec:bg_inductor}
% ............................................................

TorchInductor receives the captured FX graph and transforms it in
several stages before emitting kernel code.
The first stage is a set of \emph{post-gradient FX passes} that
rewrite subgraph patterns at the ATen operator level.
These passes perform algebraic simplifications, operator fusion
opportunities, and target-specific graph rewrites; they operate
directly on FX graph nodes and leave behind a semantically equivalent
but more execution-friendly graph.
In this work, we extend this FX pass stage to match the weight-linear
and attention sub-graphs present in each transformer layer and rewrite
them into a tensor-parallel form, inserting per-socket weight slices
and propagating kernel-sharding metadata through the graph; the details
are described in Section~\ref{sec:approach}.

Following the FX passes, Inductor lowers the transformed graph to its
own loop-level intermediate representation.  This IR makes loop bounds,
reduction axes, and access patterns explicit, enabling Inductor's
scheduler to identify kernels that can be fused, determine tiling
strategies, and allocate intermediate buffers.  The scheduler then
hands each kernel to the configured codegen backend—Triton CPU in our
setting—together with the loop structure and buffer metadata required to
generate executable code.

% ............................................................
\subsubsection{Triton Code Generation}
\label{sec:bg_triton_cg}
% ............................................................

Inductor's Triton codegen backend~\cite{ansel2024pytorch2} translates
each scheduled kernel into a Triton kernel program.
We use the existing codegen path without changes to the kernel body
representation; only the kernel \emph{call-site} wrapper is extended to
carry the per-socket tensor pointers and metadata described in
Section~\ref{sec:approach}.

In the generated kernel, computation is organised as a
three-dimensional grid of independently executable \emph{program
instances}.  Each program instance is identified at runtime by its
three integer coordinates within the grid—analogous to a thread block
index in GPU programming.  Within a program instance, operations are
applied to a contiguous tile of elements whose width is chosen to match
the SIMD vector width of the target processor.  Tile-width boundaries
and out-of-bounds access are handled through predicate masks, so the
same kernel body executes correctly regardless of whether the problem
size divides evenly into tiles.  A Python wrapper module generated
alongside the kernel encodes the grid dimensions and the argument
layout needed to invoke the compiled binary.

% ............................................................
\subsubsection{Kernel Compilation via MLIR}
\label{sec:bg_compilation}
% ............................................................

Each Triton kernel is compiled through a multi-stage MLIR lowering
pipeline~\cite{lattner2021mlir} that progressively translates the kernel
from Triton's high-level tile IR through a CPU-specific MLIR dialect—
applying vectorisation, SIMD instruction selection, and target-specific
optimisations—before emitting native assembly and linking it into a
shared object binary (\texttt{.so}).
This binary is stored in a content-addressed on-disk cache and loaded
into the Python process at inference time, so repeated invocations with
the same kernel signature and input shapes incur no recompilation
overhead.

% ............................................................
\subsubsection{Triton CPU Runtime}
\label{sec:bg_triton_rt}
% ............................................................

At runtime, each kernel binary is invoked through a
\emph{CPULauncher}: a small C++ extension module that is compiled once
per unique kernel signature and loaded dynamically into the Python
process.  The launcher's \texttt{launch} entry point receives the
kernel's three-dimensional grid extents, a pointer to the compiled
kernel function, and the tensor argument list.  It then constructs the
complete set of grid coordinates—one three-tuple for each program
instance—and distributes execution of those program instances across the
available processor cores.

In the standard execution path the launcher uses OpenMP's static-schedule
parallel loop to assign contiguous ranges of grid blocks to hardware
threads.  Each thread calls the kernel function with the coordinates of
each block in its assigned range; the kernel uses those coordinates to
compute the element range it is responsible for and performs the
corresponding vectorised load-compute-store operations.  No
synchronisation between threads is needed because grid blocks are
defined to be independent.

The current upstream Triton CPU runtime is \emph{NUMA-unaware}: threads
are allocated to grid blocks without any regard for which physical
memory controller holds the operand tensors, and the operating system
scheduler is free to migrate threads across sockets between kernel
launches.  As a consequence, a significant fraction of memory accesses
cross the inter-socket interconnect, incurring remote NUMA latency that
substantially reduces effective memory bandwidth.  Our modifications to
this runtime—which introduce a persistent per-socket thread pool with
explicit CPU affinity binding and route each socket's threads to the
weight copies resident in local DRAM—are described in
Section~\ref{sec:approach}.

% ============================================================
%  References
% ============================================================
\bibliographystyle{IEEEtran}

\begin{thebibliography}{10}

\bibitem{ansel2024pytorch2}
J.~Ansel, E.~Yang, H.~He, N.~Gimelshein, A.~Jain, M.~Voznesensky,
B.~Bao, P.~Bell, D.~Berard, E.~Burovski, G.~Chauhan, A.~Chourdia,
W.~Constable, A.~Desmaison, Z.~DeVito, E.~Ellison, W.~Feng, J.~Gong,
M.~Gschwind, B.~Hirsh, S.~Huang, K.~Kalamkar, J.~Kirsch, M.~Lazos,
M.~Lezcano, Y.~Liang, J.~Liang, Y.~Lu, C.~K.~Luk, B.~Maher,
Y.~Pan, C.~Puhrsch, M.~Reso, M.~Saroufim, M.~Y.~Siraichi,
H.~Suk, M.~Suo, P.~Tillet, E.~Wang, X.~Wang, W.~Wen, S.~Zhang,
X.~Zhao, K.~Zhou, R.~Zou, A.~Mathews, G.~Chanan, P.~Wu, and
S.~Chintala,
``PyTorch~2: Faster machine learning through dynamic Python bytecode
transformation and graph compilation,''
in \emph{Proc. 29th ACM Int. Conf. Architectural Support for
Programming Languages and Operating Systems (ASPLOS)}, vol.~2,
pp.~929--947, ACM, 2024.

\bibitem{paszke2019pytorch}
A.~Paszke, S.~Gross, F.~Massa, A.~Lerer, J.~Bradbury, G.~Chanan,
T.~Killeen, Z.~Lin, N.~Gimelshein, L.~Antiga, A.~Desmaison,
A.~Kopf, E.~Yang, Z.~DeVito, M.~Raison, A.~Tejani, S.~Chilamkurthy,
B.~Steiner, L.~Fang, J.~Bai, and S.~Chintala,
``PyTorch: An imperative style, high-performance deep learning library,''
in \emph{Advances in Neural Information Processing Systems},
vol.~32, pp.~8026--8037, 2019.

\bibitem{tillet2019triton}
P.~Tillet, H.~T.~Kung, and D.~Cox,
``Triton: An intermediate language and compiler for tiled neural network
computations,''
in \emph{Proc. 3rd ACM SIGPLAN Int. Workshop on Machine Learning and
Programming Languages (MAPL)}, pp.~10--19, ACM, 2019.

\bibitem{triton_cpu_repo}
OpenAI Triton Contributors,
``Triton-CPU: CPU backend for the Triton compiler,''
GitHub repository.
[Online]. Available: \url{https://github.com/triton-lang/triton-cpu}

\bibitem{lattner2021mlir}
C.~Lattner, M.~Amini, U.~Bondhugula, A.~Cohen, A.~Davis, J.~Pienaar,
R.~Riddle, T.~Shpeisman, N.~Vasilache, and O.~Zinenko,
``MLIR: Scaling compiler infrastructure for domain specific
computation,''
in \emph{Proc. IEEE/ACM Int. Symp. Code Generation and Optimization
(CGO)}, pp.~2--14, IEEE, 2021.

\bibitem{vaswani2017attention}
A.~Vaswani, N.~Shazeer, N.~Parmar, J.~Uszkoreit, L.~Jones,
A.~N.~Gomez, \L.~Kaiser, and I.~Polosukhin,
``Attention is all you need,''
in \emph{Advances in Neural Information Processing Systems},
vol.~30, 2017.

\bibitem{brown2020gpt3}
T.~B.~Brown, B.~Mann, N.~Ryder, M.~Subbiah, J.~Kaplan, P.~Dhariwal,
A.~Neelakantan, P.~Shyam, G.~Sastry, A.~Askell \emph{et~al.},
``Language models are few-shot learners,''
in \emph{Advances in Neural Information Processing Systems},
vol.~33, pp.~1877--1901, 2020.

\bibitem{radford2019gpt2}
A.~Radford, J.~Wu, R.~Child, D.~Luan, D.~Amodei, and I.~Sutskever,
``Language models are unsupervised multitask learners,''
\emph{OpenAI Blog}, vol.~1, no.~8, p.~9, 2019.

\bibitem{shoeybi2019megatron}
M.~Shoeybi, M.~Narayanan, H.~Peng, C.~Chi, J.~Casper, and
B.~Catanzaro,
``Megatron-LM: Training multi-billion parameter language models using
model parallelism,''
\emph{arXiv preprint arXiv:1909.08053}, 2019.

\bibitem{luk2005intel}
C.~K.~Luk, R.~Cohn, R.~Muth, H.~Patil, A.~Klauser, G.~Lowney,
S.~Wallace, V.~J.~Reddi, and K.~Hazelwood,
``Pin: Building customized program analysis tools with dynamic
instrumentation,''
in \emph{Proc. ACM SIGPLAN Conf. Programming Language Design and
Implementation (PLDI)}, pp.~190--200, ACM, 2005.

\end{thebibliography}

\end{document}
